% TODO make hw stylesheet
\documentclass[10pt]{article}
\usepackage[letterpaper,top=1in]{geometry}
\usepackage{quiver}
\usepackage[dvipsnames,svgnames]{xcolor}
\usepackage[shortlabels]{enumitem}
\usepackage[framemethod=TikZ]{mdframed}
\usepackage{amsmath,amssymb,amsthm}
\usepackage[colorlinks]{hyperref}
\usepackage{multicol}
\usepackage{tabularx}
\usepackage{stmaryrd}

\hypersetup{
	colorlinks=true,
	linkcolor=blue,
	filecolor=magenta,
	urlcolor=magenta,
	pdftitle={MAT 534 HW}
}

\usepackage{mathtools}
\usepackage{thmtools}
\usepackage[textsize=footnotesize]{todonotes}
\usepackage{titling}

\reversemarginpar
\mdfdefinestyle{mdbluebox}{roundcorner=10pt,innerbottommargin=9pt,
	linecolor=blue,backgroundcolor=TealBlue!5,}
\declaretheoremstyle[headfont=\sffamily\bfseries\color{MidnightBlue},
mdframed={style=mdbluebox},]{thmbluebox}

\mdfdefinestyle{mdbloobox}{frametitlefont=\bfseries,innerbottommargin=8pt,
	nobreak=true,backgroundcolor=TealBlue!5,linecolor=blue,}
\declaretheoremstyle[headfont=\bfseries\color{MidnightBlue},
mdframed={style=mdbloobox},headpunct={\\[3pt]},postheadspace=0pt,]{thmbloobox}

\mdfdefinestyle{mdredbox}{frametitlefont=\bfseries,innerbottommargin=8pt,
	nobreak=true,backgroundcolor=Salmon!5,linecolor=RawSienna,}
\declaretheoremstyle[headfont=\bfseries\color{RawSienna},
mdframed={style=mdredbox},headpunct={\\[3pt]},postheadspace=0pt,]{thmredbox}

\mdfdefinestyle{mdgreenbox}{linecolor=ForestGreen,backgroundcolor=ForestGreen!5,
	linewidth=2pt,rightline=false,leftline=true,topline=false,bottomline=false,}
\declaretheoremstyle[headfont=\bfseries\sffamily\color{ForestGreen!70!black},
mdframed={style=mdgreenbox},headpunct={ --- },]{thmgreenbox}

\mdfdefinestyle{mdorangebox}{linecolor=RedOrange,backgroundcolor=RedOrange!5,
	linewidth=2pt,rightline=false,leftline=true,topline=false,bottomline=false,}
\declaretheoremstyle[headfont=\bfseries\sffamily\color{RedOrange!70!black},
mdframed={style=mdorangebox},headpunct={ --- },]{thmorangebox}

\mdfdefinestyle{mdblackbox}{linecolor=black,backgroundcolor=RedViolet!5!gray!5,
	linewidth=3pt,rightline=false,leftline=true,topline=false,bottomline=false,}
\declaretheoremstyle[mdframed={style=mdblackbox}]{thmblackbox}

\declaretheoremstyle[spaceabove=6pt,spacebelow=6pt]{basehead}

\declaretheorem[style=basehead,name=Problem]{problem}
\declaretheorem[style=thmredbox,name=Problem,numbered=no]{problem*}
\declaretheorem[style=thmbloobox,name=Setup]{setup} % for config geo
\declaretheorem[style=thmredbox,name=Example,sibling=problem]{example}
\declaretheorem[style=thmredbox,name=$\bigstar$ Problem,sibling=problem]{niceprob}
\declaretheorem[style=thmbluebox,name=Theorem,sibling=problem]{theorem}
\declaretheorem[style=thmbluebox,name=Corollary,sibling=theorem]{corollary}
\declaretheorem[style=thmbluebox,name=Proposition,sibling=theorem]{prop}
\declaretheorem[style=thmbluebox,name=Lemma,numberwithin=problem]{lemma}
\declaretheorem[style=thmbluebox,name=Theorem,numbered=no]{theorem*}
\declaretheorem[style=thmbluebox,name=Lemma,numbered=no]{lemma*}
\declaretheorem[style=thmgreenbox,name=Claim,numberwithin=problem]{claim}
\declaretheorem[style=thmgreenbox,name=Claim,numbered=no]{claim*}
\declaretheorem[style=thmblackbox,name=Remark,numberwithin=problem]{remark}
\declaretheorem[style=thmblackbox,name=Remark,numbered=no]{remark*}
\declaretheorem[style=thmgreenbox,name=Definition,sibling=theorem]{definition}
\declaretheorem[style=thmgreenbox,name=Definition,numbered=no]{definition*}

\providecommand{\ol}{\overline}
\providecommand{\quiv}{\equiv}
\providecommand{\ceil}[1]{\left\lceil #1 \right\rceil}
\providecommand{\floor}[1]{\left\lfloor #1 \right\rfloor}
\newcommand{\dd}{\mathrm{d}}
\providecommand{\ul}{\underline}
\providecommand{\eps}{\varepsilon}
\providecommand{\half}{\frac{1}{2}}
\providecommand{\CC}{\mathbb C}
\providecommand{\FF}{\mathbb F}
\providecommand{\NN}{\mathbb N}
\providecommand{\QQ}{\mathbb Q}
\providecommand{\RR}{\mathbb R}
\providecommand{\ZZ}{\mathbb Z}
\providecommand{\dg}{^\circ}
\providecommand{\ii}{\item}
\providecommand{\setdiff}{\smallsetminus}
\providecommand{\alert}[1]{{\sffamily\textbf{\textcolor{blue}{#1}}}}
\providecommand{\inv}{^{-1}}
\providecommand{\mb}[1]{\mathbf{#1}}

\newenvironment{soln}{\begin{proof}[Solution]}{\end{proof}}

\providecommand{\pow}{\operatorname{Pow}}
\providecommand{\aut}{\operatorname{Aut}}
\providecommand{\vocab}[1]{{\textbf{\textcolor{ForestGreen}{#1}}}}
\providecommand{\lcm}{\operatorname{lcm}}
\providecommand{\abs}[1]{\left|#1\right|}
\providecommand{\Ob}{\operatorname{Ob}}
\providecommand{\Hom}{\operatorname{Hom}}
\providecommand{\Aut}{\operatorname{Aut}}
\providecommand{\id}{\operatorname{id}}
\newcommand{\doubleparen}[1]{(\!(#1)\!)}

\usepackage{mathrsfs}

% place title at top right
\setlength{\droptitle}{-8ex}
\pretitle{\begin{flushright}\Large\bfseries}
\posttitle{\par\end{flushright}}
\preauthor{\begin{flushright}}
\postauthor{\end{flushright}}
\predate{\begin{flushright}}
\postdate{\end{flushright}}

\title{MAT 534 HW06}
\author{\large Aditya Pahuja}
\date{\large Due 10/22}
\begin{document}
\maketitle

\begin{problem}
    Let $R$ be a ring with $1$.
    Show that if $0=1$, then $R$ is the zero ring.
\end{problem}

\begin{soln}
    Let $x\in R$. Then, using the result of Problem 2a,
    \begin{equation*}
        x = 1\cdot x = 0\cdot x = 0
    \end{equation*}
    so $R$ has exactly one element and is the zero ring.
\end{soln}

\begin{problem}
    Let $R$ be a ring. We use $-a$ to denote the additive inverse of $a\in R$.
    \begin{enumerate}[(a)]
        \ii Prove that $0\cdot a = 0$ for every $a\in R$.
	\ii Prove that $(-a)(-b) = ab$ for every $a,b\in R$.
    \end{enumerate}
\end{problem}

\begin{soln}
    (a) Since $0 + 0 = 0$,
    \begin{equation*}
        0\cdot a = (0 + 0)\cdot a = 0\cdot a + 0\cdot a,
    \end{equation*}
    and since $R$ is a group under $+$, we can cancel $0\cdot a$ from both sides
    to get $0 = 0\cdot a$.
    \\

    (b) First, we compute
    \begin{align*}
        0 = 0\cdot y = (x + (-x))\cdot y = xy + (-x)y
	&\implies -(xy) = (-x)y
	\tag{1}\\
        0 = x\cdot 0 = x\cdot(y + (-y)) = xy + x(-y)
	&\implies -(xy) = x(-y)
	\tag{2}
    \end{align*}
    (where $0 = x\cdot (0 + 0) - x\cdot 0 = x\cdot 0$)
    for any $x,y\in R$. Then
    \begin{equation*}
	ab = -(-(ab)) \overset{(1)}= -((-a)b) \overset{(2)}= (-a)(-b).\qedhere
    \end{equation*}
\end{soln}

\begin{problem}
    Prove that every ideal in the ring $\ZZ$ is of the form $(n)$
    for some $n\in\ZZ$.
\end{problem}

\begin{soln}
    Given ideal $I\subseteq\ZZ$, let $n$ be the smallest positive integer in $I$.
    Then, for any $m\in I$, we can apply the division algorithm to get that
    $m = nq + r$ for some integers $q$ and $r$ such that $0 \le r < n$.
    Since $I$ is a group under addition,
    $m - nq = r$ is in $I$, so $0\le r < n$ implies that
    $r$ must be zero, since we would otherwise contradict the minimality of $n$;
    this means $m = nq$ is a multiple of $n$.
    That is, every element of $I$ is a multiple of $n$,
    and conversely $n\in I$ means $nk\in I$ for all $k\in\ZZ$, so $I = (n)$.
\end{soln}

\begin{problem}
    Let $F$ be a field, and let $R$ be the set of all functions
    $a\colon\ZZ_{\ge0}\to F$. Here we think of $a\in R$ as representing
    the formal power series
    \begin{equation*}
	\sum_{n=0}^\infty a(n)x^n = a(0) + a(1)x + a(2)x^2 + \cdots
    \end{equation*}
    The ring of formal power series is usually denoted by $F\llbracket x\rrbracket$.
    \begin{enumerate}[(a)]
        \ii Show that $R$ is a commutative ring with $1$, under the operations
	\begin{equation*}
	    (a+b)(n) = a(n) + b(n)\quad\text{and}\quad
	    (a\cdot b)(n) = \sum_{i=0}^n a(i)b(n-i).
	\end{equation*}
	\ii Determine the set of units $R^\times$.
	\ii Determine the ideals in the ring $R$.
    \end{enumerate}
\end{problem}

\begin{soln}
    (a) The addition operation inherits associativity from the addition of $F$;
    the zero element is the function that is constantly zero;
    and the inverse of $a\colon\ZZ_{\ge0}\to F$ is
    $-a\colon\ZZ_{\ge0}\to F$ defined by $(-a)(n) = -(a(n))$.

    The multiplication operation can be written more symmetrically as
    \begin{equation*}
	(a\cdot b)(n) = \sum_{i+j = n}a(i)b(j),
    \end{equation*}
    which lets us recover commutativity immediately because $a(i)b(j) = b(j)a(i)$
    and associativity because
    \begin{align*}
	((a\cdot b)\cdot c)(n)
	= \sum_{m+k = n}\sum_{i+j=m}a(i)b(j)c(k)
	&= \sum_{i+j+k=n}a(i)b(j)c(k)\\
	&= \sum_{i+m'=n}\sum_{j+k=m'}a(i)b(j)c(k)
	= (a\cdot(b\cdot c))(n).
    \end{align*}
    The multiplicative identity in this ring
    is the function $e\colon\ZZ_{\ge0}\to F$ given by
    $e(0) = 1$ and $e(n) = 0$ for $n\ge 1$, since
    \begin{equation*}
	(a\cdot e)(n) = \sum_{i=0}^n a(n-i)e(i)
	= a(n)e(0) + \sum_{i=1}^n a(n-i)\underbrace{e(i)}_{0} = a(n).
    \end{equation*}

    (b) The set of units of $R$ is the set of elements $a$ with
    $a(0)\ne 0$. To check that all such elements are invertible,
    we can construct an inverse $b$ as such:
    \begin{itemize}
        \ii Set $b(0) = a(0)\inv$, so $(a\cdot b)(0) = 1$.
	\ii For $n\ge 1$, we want
	\begin{equation*}
	    (a\cdot b)(n) = \sum_{i=0}^n a(n-i)b(i) = 0,
	\end{equation*}
	so, solving for $b(n)$ here, we can choose
	\begin{equation*}
	    b(n) \coloneq a(0)\inv\sum_{i=0}^{n-1} a(n-i)b(i)
	\end{equation*}
	to force this, given $b(n-1)$, \dots, $b(0)$.
    \end{itemize}
    On the other hand, if $a(0) = 0$, then for any $b\in R$
    $(a\cdot b)(0) = a(0)b(0) = 0\ne 1$, so $a$ cannot be invertible.
\end{soln}

\begin{problem}
    Find a ring $R_0$ in \doubleparen{unital commutative rings}
    such that, for any unital commutative ring $R$,
    the set $\Hom(R_0,R)$ is in one-to-one correspondence with $R$.
\end{problem}

\begin{soln}
    Let $R_0 = \ZZ[x]$. Then a morphism $\phi\colon R_0\to R$
    is uniquely determined by $\phi(x)$ (since $\phi(1) = 1$ is fixed),
    and $\phi(x)$ can be set to any element of $R$,
    so there is a morphism for every element of $R$.
\end{soln}

\begin{problem}
    Find a ring $R_1$ in \doubleparen{unital commutative rings}
    such that, for any unital commutative ring $R$,
    the set $\Hom(R_1,R)$ is in one-to-one correspondence with $R^\times$.
\end{problem}

\begin{soln}
    Let $R_1 = \ZZ[x,1/x]$ be the ring of Laurent polynomials over $\ZZ$.
    Then a morphism $\phi\colon R_1\to R$
    is uniquely determined by $\phi(x)$ (since $\phi(1) = 1$ is fixed).
    Let $\phi_r$ be the morphism sending $x$ to $r\in R$.
    On one hand $r$ must be a unit for $\phi_r$ to be well-defined
    because $\phi_r(x)\phi_r(x\inv) = \phi_r(1) = 1$
    so $\phi_r(x) = r$ has an inverse; on the other hand,
    if $r$ is a unit then $\phi_r$ is well-defined,
    being the evaluation morphism $p(x)\mapsto p(r)$;
    thus $r\mapsto \phi_r$ is a bijection from $R$ to $\Hom(R_1,R)$.
\end{soln}

\begin{problem}
    If $R$ is a unital ring, the set
    $C(R) = \{x\in R : xy = yx \text{ for all }y\in R\}$
    is called the \vocab{center} of $R$.
    Elements of $C(R)$ are called \vocab{central}.
    \begin{enumerate}[(a)]
        \ii Show that $C(R)$ is a subring of $R$.
	\ii Let $R = M_n(F)$ be the ring of $n\times n$ matrices over field $F$.
	What is the center of $R$?
	\ii Let $G$ be a finite group, and let $R$ be
	the integral group ring $\ZZ[G]$ of $G$.
	What is the center of $R$?
    \end{enumerate}
\end{problem}

\begin{soln}
    (a) Clearly $1$ is central. If $a$ and $b$ are central then
    for all $x\in R$, $(a + b)x = ax + bx = xa + xb = x(a + b)$
    so $a + b$ is central, and $abx = axb = xab$ implies $ab$ is central.
    These three statements imply $C(R)$ is a subring of $R$.
    \\

    (b) The center of $R$ is the diagonal matrices $\lambda I$, $\lambda\in F$
    (where $I$ is the identity matrix). Clearly these commute with all elements of $R$,
    so we need to show that any central element is of this form.

    Let $E_{ij}\in R$ be the matrix whose entries are all zero
    except for a $1$ in row $i$, column $j$. For any $A\in R$,
    $E_{ij}A$ is the matrix whose $i$th row is equal to the $j$th row of $A$
    and all other entries are zero,
    while $AE_{ij}$ is the matrix whose $j$th column is equal to
    the $i$th column of $A$ and all other entries are zero.

    If $A$ is central in $R$, then
    these two matrices are equal for all $i,j\in\{1,\dots,n\}$.
    The entry at row $i$, column $j$ in the product is equal to both
    the entry at row $j$, column $j$ and the entry at row $i$, column $i$
    in $A$, meaning as $i$ and $j$ vary we see that
    $A$ has the same element of $F$ at every entry in the diagonal.
    Moreover, if $i\ne j$ then the entry of $E_{ji}A$
    at row $j$, column $j$ is equal to the row $i$, column $j$ entry of $A$,
    and the corresponding entry of $AE_{ji}$ is equal to zero
    since the nonzero row is row $i$; therefore the entry in $A$
    at row $i$, column $j$ is zero, so any entry that is
    off the main diagonal must be zero, i.e. $A$ is a multiple of $I$.
    \\

    (c) Let $C_1$, $C_2$, \dots, $C_n$ be the conjugacy classes of $G$.
    Then let $s_j = \sum_{g\in C_j} g$ and let
    \begin{equation*}
	Z = \left\{\sum_{j=1}^n c_js_j : c_1,\dots,c_n\in\ZZ\right\}.
    \end{equation*}
    It is easy to see that $Z\subseteq C(\ZZ[G])$, since
    if $h\in G$ then $hC_j = C_jh$, so that $hs_j = s_jh$.
    Conversely, if $\sum_{x\in G} c_xx$ is central, then, conjugating by $g$,
    \begin{equation*}
	\sum_{x\in G} c_xx
	= \sum_{x\in G} c_x gxg\inv
	= \sum_{x\in G} c_{g\inv xg} x,
    \end{equation*}
    so if $y$ is conjugate to $x$ then $c_y = c_x$, i.e.
    the coefficients of any central element
    are constant on conjugacy classes,
    so the central element is in $Z$.
\end{soln}

\begin{problem}
    Suppose that $f(x,y)$ is a polynomial in $\ZZ[x,y]$, and
    $R = \ZZ[x,y]/(f(x,y))$. Show that $\Hom(R,\RR)$ is in
    one-to-one correspondence with the points on the graph of
    $f(x,y) = 0$ in $\RR^2$.
\end{problem}

\begin{soln}
    $\Hom(R,\RR)$ is in one-to-one correspondence with
    the maps $\phi\colon\ZZ[x,y]\to\RR$ such that $(f(x,y))\subseteq\ker\phi$.
    Let $\phi_{r,s}$ be the ring homomorphism
    $\ZZ[x,y]\to\RR$ sending $x$ to $r$ and $y$ to $s$;
    the kernel of this map contains $(f)$ if and only if
    $\phi_{r,s}(f(x,y)) = f(\phi_{r,s}(x), \phi_{r,s}(y)) = f(r,s) = 0$,
    which gives the bijection between points $(r,s) \in f\inv(0)$
    and morphisms $\phi_{r,s}\in\Hom(R,\RR)$.
\end{soln}

\begin{problem}
    Let $R$ be a unital ring.
    \begin{enumerate}[(a)]
        \ii If $e\in R$ is a central idempotent, show that
	$R\cong R/(e)\times R/(1-e)$.
	\ii Let $H$ be a subgroup of a finite group $G$.
	Show that the element $e_H = \frac 1{\abs H}\sum_{g\in H}g$
	is an idempotent in the group ring $\QQ[G]$.
	Under what conditions on $H$ does $e_H$ belong to
	the center of the group ring?
    \end{enumerate}
\end{problem}

\begin{soln}
    (a) If $e$ is a central idempotent then so is $1-e$.
    Clearly $(e) + (1-e) = (1)$
    and the intersection $(e)\cap(1-e)$ is trivial since
    if $ae = b(1-e)$ for some $a,b\in R$ then this common element is
    $ae = ae^2 = b(1-e)e = 0$.
    The Chinese remainder theorem then implies that
    $R\cong R/(e)\times R/(1-e)$.
    \\
    
    (b) We can compute
    \begin{equation*}
	e^2_H = \frac 1{|H|^2} \sum_{x\in H}\sum_{g\in H} xg
	= \frac1{|H|^2} \cdot \abs H\cdot\sum_{g\in H}g
	= \frac 1{\abs H}\sum_{g\in H}g = e_H
    \end{equation*}
    since $g\mapsto xg$ is a bijection on $H$.
    By the same logic to 7(c) (the group ring being over $\QQ$
    and not $\ZZ$ changes nothing), $e_H$ is central if and only if
    $H$ is a union of conjugacy classes of $G$, i.e. $H$ is normal.
\end{soln}

\begin{problem}
    Let $X$ be a compact Hausdorff space, and let $C(X)$
    be the ring of continuous functions $f\colon X\to\RR$.
    \begin{enumerate}[(a)]
        \ii Is $C(X)$ an integral domain?
	\ii Let $x\in X$ be any point. Show that
	$\{f\in C(X) : f(x) = 0\}$ is a maximal ideal of $C(X)$.
	\ii Show that every maximal ideal of $C(X)$ is of this form.
    \end{enumerate}
\end{problem}

\begin{soln}
    (a) No; let $A\subsetneq X$ be closed with nonempty interior
    and $B = \ol{X\setminus A} \ne X$ also be closed, so $A\cup B = X$.
    Then, $C(X)$ contains a nonzero function that is zero on $A$
    (concrete construction:
    pick a point $p\notin A$ and, since compact Hausdorff spaces are normal,
    use Urysohn's lemma to get a function that's zero on $A$ and $1$ at $p$)
    and a nonzero function that is zero on $B$.
    The product of these two functions is zero on $A\cup B$,
    so the two functions are zero divisors in $C(X)$.
    \\

    (b) Let $M_x = \{f\in C(X) : f(x) = 0\}$,
    and let $g\in C(X)$ be a function that is not zero at $x$.
    Then, $B = g\inv(0)\ne X$ is a closed subset of $X$.
    Since $X$ is compact Hausdorff, it is normal, so
    there exists, by Urysohn's lemma, a continuous function
    $f\colon X\to\RR$ such that $f|_A = 0$ and $f(p) = 1$
    where $p$ is a point in $X\setminus A$.
    Then, $f^2 + g^2$ is a continuous function that is nowhere zero,
    which means it is a unit in $C(X)$, as it has inverse
    $\frac{1}{f^2 + g^2}$. That is, the ideal generated by $M_x$
    and any $g$ outside $M_x$ is equal to the entire ring $C(X)$,
    so $M_x$ must be maximal.
    \\

    (c) Let $M$ be an ideal.
    If, for each $x\in X$, $M$ contains a function that is nonzero at $x$,
    then the sets $f\inv(\RR\setminus\{0\})$ are an open cover of $X$,
    so we can choose finitely many of these functions $f_1,\dots,f_n$
    where $f_1\inv(\RR\setminus\{0\})\cup\cdots\cup f_n\inv(\RR\setminus\{0\}) = X$.
    This means $f_1^2 + \cdots + f_n^2\in M$ is nowhere zero, so $M$
    contains a unit, hence $M = C(X)$.

    Thus, if $M$ is maximal, there exists a closed subset $A\subseteq X$ such that
    $f|_A = 0$ for all $f\in M$. However this ideal is contained in $M_x$
    for any $x\in A$, so the $M_x$ are the only such ideals that are maximal.
\end{soln}

\begin{problem}
    Let $R$ be a unital commutative ring.
    Show that an element $r\in R$ is a unit if and only if
    $r$ is not contained in any maximal ideal of $R$.
\end{problem}

\begin{soln}
    If $r$ is a unit then an ideal containing $r$
    is the entire ring, so it cannot be maximal.

    If $r$ is not a unit then consider the set of all
    proper ideals of $R$ that contain $r$, ordered by inclusion.
    Given a totally ordered subcollection of these ideals $\{I_\alpha\}_\alpha$,
    the union $\bigcup_\alpha I_\alpha$ is also an ideal,
    and this ideal is not $R$ because none of the $I_\alpha$ contains a unit,
    so it is an upper bound for the totally ordered subcollection.
    By Zorn's lemma, the collection of proper ideals containing $r$
    has a maximal element, which is a maximal ideal containing $r$.
\end{soln}











\end{document}
